\chapter{2022年浙江卷}

\section{单选题}

\begin{problem}
设集合 \(A=\{1,2\}\)，\(B=\{2,4,6\}\)，则 \(A\cup B=\)（\quad）

\choices
    {\(\{2\}\)}
    {\(\{1,2\}\)}
    {\(\{2,4,6\}\)}
    {\(\{1,2,4,6\}\)}

\end{problem}

\begin{answer}
D
\end{answer}

\begin{solution}
因为 \(A=\{1,2\}\)，\(B=\{2,4,6\}\)，
\(A\cup B=\{1,2,4,6\}\)，
故选：D.
\end{solution}
\begin{problem}
已知 \(a,b\in\R\)，且 \(a+3\i=(b+\i)\i\)，则（\quad）

\choices
    {\(a=1\)，\(b=-3\)}
    {\(a=-1\)，\(b=3\)}
    {\(a=-1\)，\(b=-3\)}
    {\(a=1\)，\(b=3\)}

\end{problem}

\begin{answer}
B
\end{answer}

\begin{solution}
因为
\(a+3\i=(b+\i)\i=-1+b\i\)，
且 \(a,b\in \R\)，
\(a=-1\)，\(b=3\)，
故选：B.
\end{solution}
\begin{problem}
若实数 \(x,y\) 满足约束条件 \(x-2\ge 0\)，\(2x+y-7\le 0\)，\(x-y-2\le 0\)
则 \(z=3x+4y\) 的最大值是（\quad）

\choices
    {20}
    {18}
    {13}
    {6}

\end{problem}

\begin{answer}
B
\end{answer}

\begin{solution}
实数 \(x,y\) 满足约束条件
\[
\begin{cases}
x-2\ge 0,\\
2x+y-7\le 0,\\
x-y-2\le 0,
\end{cases}
\]
则不等式组表示的平面区域为如图所示的阴影部分.

\begin{center}
\bitmapfigure[max width=0.28\textwidth,max totalheight=4.8cm]{img/2022/zhejiang/q3_solution_fig1.png}
\end{center}

由已知可得 \(A(2,3)\)，由图可知：当直线 \(3x+4y-z=0\) 过点 \(A\) 时，\(z\) 取最大值，
\(z=3x+4y=3\times 2+4\times 3=18\)，
故选：B.
\end{solution}
\begin{problem}
设 \(x\in\R\)，则“\(\sin x=1\)”是“\(\cos x=0\)”的（\quad）

\choices
    {充分不必要条件}
    {必要不充分条件}
    {充分必要条件}
    {既不充分也不必要条件}

\end{problem}

\begin{answer}
A
\end{answer}

\begin{solution}
因为 \(\sin^2 x+\cos^2 x=1\)，

\circled{1}当 \(\sin x=1\) 时，则 \(\cos x=0\)，所以充分性成立.

\circled{2}当 \(\cos x=0\) 时，则 \(\sin x=\pm 1\)，所以必要性不成立.

所以“\(\sin x=1\)”是“\(\cos x=0\)”的充分不必要条件，故选：A.
\end{solution}
\begin{problem}
某几何体的三视图如图所示（单位：cm），则该几何体的体积（单位：cm\(^3\)）是（\quad）

\begin{center}
\bitmapfigure[max width=0.34\textwidth,max totalheight=4.8cm]{img/2022/zhejiang/q5_fig1.png}
\end{center}

\choices
    {\(22\pi\)}
    {\(8\pi\)}
    {\(\dfrac{22\pi}{3}\)}
    {\(\dfrac{16\pi}{3}\)}

\end{problem}

\begin{answer}
C
\end{answer}

\begin{solution}
由三视图可知几何体是上部为半球，中部是圆柱，下部是圆台，

所以几何体的体积为
\(\frac12\times \frac{4\pi}{3}\times 1^3+\pi\times 1^2\times 2 +\frac13\left(2^2\times \pi+1^2\times \pi+\sqrt{2^2\times \pi\times 1^2\times \pi}\right)\times 2 =\frac{22\pi}{3}\).
故选：C.
\end{solution}
\begin{problem}
为了得到函数 \(y=2\sin 3x\) 的图象，只要把函数 \(y=2\sin\left(3x+\frac{\pi}{5}\right)\) 图象上所有的点（\quad）

\choices
    {向左平移 \(\dfrac{\pi}{5}\) 个单位长度}
    {向右平移 \(\dfrac{\pi}{5}\) 个单位长度}
    {向左平移 \(\dfrac{\pi}{15}\) 个单位长度}
    {向右平移 \(\dfrac{\pi}{15}\) 个单位长度}

\end{problem}

\begin{answer}
D
\end{answer}

\begin{solution}
因为
\(2\sin 3x=2\sin\left(3\left(x-\frac{\pi}{15}\right)+\frac{\pi}{5}\right)\)，
所以把函数 \(y=2\sin\left(3x+\frac{\pi}{5}\right)\) 图象上的所有点向右平移 \(\frac{\pi}{15}\) 个单位长度即可得到函数 \(y=2\sin 3x\) 的图象.

故选：D.
\end{solution}
\begin{problem}
已知 \(2^a=5\)，\(\log_8 3=b\)，则 \(4^{a-3b}=\)（\quad）

\choices
    {25}
    {5}
    {\(\dfrac{25}{9}\)}
    {\(\dfrac{5}{3}\)}

\end{problem}

\begin{answer}
C
\end{answer}

\begin{solution}
由 \(2^a=5\)，\(\log_8 3=b\)，
\(8^b=3\)，\(2^{3b}=3\).
则
\(4^{a-3b}=\frac{4^a}{4^{3b}}=\frac{(2^a)^2}{(2^{3b})^2}=\frac{5^2}{3^2}=\frac{25}{9}\).
故选：C.
\end{solution}
\begin{problem}
如图，已知正三棱柱 \(ABC-A_1B_1C_1\)，\(AC=AA_1\)，\(E,F\) 分别是棱 \(BC,A_1C_1\) 上的点. 记 \(EF\) 与 \(AA_1\) 所成的角为 \(\alpha\)，\(EF\) 与平面 \(ABC\) 所成的角为 \(\beta\)，二面角 \(F-BC-A\) 的平面角为 \(\gamma\)，则（\quad）

\begin{center}
\bitmapfigure[max width=0.28\textwidth,max totalheight=4.8cm]{img/2022/zhejiang/q8_fig1.png}
\end{center}

\choices
    {\(\alpha\le \beta\le \gamma\)}
    {\(\beta\le \alpha\le \gamma\)}
    {\(\beta\le \gamma\le \alpha\)}
    {\(\alpha\le \gamma\le \beta\)}

\end{problem}

\begin{answer}
A
\end{answer}

\begin{solution}
\begin{center}
\bitmapfigure[max width=0.28\textwidth,max totalheight=4.8cm]{img/2022/zhejiang/q8_solution_fig1.png}
\end{center}

因为正三棱柱 \(ABC-A_1B_1C_1\) 中，\(AC=AA_1\)，所以正三棱柱的所有棱长相等，设棱长为 1.

如图，过 \(F\) 作 \(FG\perp AC\)，垂足点为 \(G\)，连接 \(GE\)，则 \(A_1A\parallel FG\)，
\(EF \text{ 与 } AA_1 \text{ 所成的角为 } \angle EFG=\alpha\)，
且
\(\tan\alpha=\frac{GE}{FG}=GE\).

又 \(GE\in [0,1]\)，所以 \(\tan\alpha\in [0,1]\).

因为 \(EF\) 与平面 \(ABC\) 所成的角为 \(\angle FEG=\beta\)，且
\(\tan\beta=\frac{FG}{GE}=\frac{1}{GE}\).
所以 \(\tan\beta\in [1,+\infty)\)，从而 \(\tan\beta\ge \tan\alpha\).

再过 \(G\) 点作 \(GH\perp BC\)，垂足点为 \(H\)，连接 \(HF\).

又因为 \(A_1A\perp\) 底面 \(ABC\)，且 \(A_1A\parallel FG\)，所以 \(FG\perp\) 底面 \(ABC\). 又 \(BC\subset\) 底面 \(ABC\)，所以 \(BC\perp FG\). 因为 \(GH\perp BC\)，且 \(FG\cap GH=G\)，所以 \(BC\perp\) 平面 \(GHF\).

因此二面角 \(F-BC-A\) 的平面角为 \(\angle GHF=\gamma\)，且
\(\tan\gamma=\frac{FG}{GH}=\frac{1}{GH}\).
又 \(GH\in \left[0,\frac{\sqrt3}{2}\right]\)，
\(\tan\gamma\in \left[\frac{2\sqrt3}{3},+\infty\right)\)，
所以 \(\tan\gamma\ge \tan\alpha\).

又 \(GE\ge GH\)，所以 \(\tan\beta\le \tan\gamma\).

\(\tan\alpha\le \tan\beta\le \tan\gamma\).
又 \(\alpha,\beta,\gamma\in \left[0,\frac{\pi}{2}\right)\)，\(y=\tan x\) 在 \(\left[0,\frac{\pi}{2}\right)\) 单调递增，
\(\alpha\le \beta\le \gamma\).
故选：A.
\end{solution}
\begin{problem}
已知 \(a,b\in\R\). 若对任意 \(x\in\R\)，\(a|x-b|+|x-4|-|2x-5|\ge 0\)，
则（\quad）

\choices
    {\(a\le 1\)，\(b\ge 3\)}
    {\(a\le 1\)，\(b\le 3\)}
    {\(a\ge 1\)，\(b\ge 3\)}
    {\(a\ge 1\)，\(b\le 3\)}

\end{problem}

\begin{answer}
D
\end{answer}

\begin{solution}
由题意有：对任意的 \(x\in \R\)，有
\(a|x-b|\ge |2x-5|-|x-4|\)
恒成立.

设
\(f(x)=a|x-b|\)，
\[
g(x)=|2x-5|-|x-4|=\begin{cases}
1-x,&x\le \frac52,\\
3x-9,&\frac52<x<4,\\
x-1,&x\ge 4.
\end{cases}
\]
即 \(f(x)\) 的图象恒在 \(g(x)\) 的上方（可重合），如下图所示：

\begin{center}
\bitmapfigure[max width=0.30\textwidth,max totalheight=4.8cm]{img/2022/zhejiang/q9_solution_fig1.png}
\end{center}

由图可知，\(a\ge 3\)，\(1\le b\le 3\)，或 \(1\le a<3\)，\(1\le b\le 4-\frac{3}{a}\le 3\).

故选：D.
\end{solution}
\begin{problem}
已知数列 \(\{a_n\}\) 满足 \(a_1=1\)，\(a_{n+1}=a_n-\frac13 a_n^2\)（\(n\in\N^*\)），
则（\quad）

\choices
    {\(2<100a_{100}<\dfrac52\)}
    {\(\dfrac52<100a_{100}<3\)}
    {\(3<100a_{100}<\dfrac72\)}
    {\(\dfrac72<100a_{100}<4\)}

\end{problem}

\begin{answer}
B
\end{answer}

\begin{solution}
因为
\(a_{n+1}-a_n=-\frac13a_n^2<0\)，
所以 \(\{a_n\}\) 为递减数列.

又 \(a_1=1\)，\(\{a_n\}\) 为递减数列，所以 \(a_n\le 1\). 且 \(a_n\ne 0\)，所以
\(\frac{a_{n+1}}{a_n}=1-\frac13a_n\ge \frac23>0\).
又 \(a_1=1>0\)，则 \(a_n>0\).

所以
\(a_n-a_{n+1}=\frac13a_n^2\ge \frac13a_na_{n+1}\)，
从而
\(\frac1{a_{n+1}}-\frac1{a_n}\ge \frac13\).
累加可得
\(\frac1{a_n}\ge \frac1{a_1}+\frac13(n-1)=\frac13n+\frac23\)，
则
\(a_n\le \frac{3}{n+2}\).
所以
\(100a_{100}\le 100\times \frac{3}{102}<\frac{306}{102}=3\).

由 \(a_{n+1}=a_n-\frac13a_n^2\) 得
\(a_{n+1}=a_n\left(1-\frac13a_n\right)\)，
得
\(\frac1{a_{n+1}}-\frac1{a_n} =\frac{1}{3-a_n} \le \frac{1}{3-\frac{3}{n+2}} =\frac13\left(1+\frac1{n+1}\right)\)，
累加可得
\(\frac1{a_n}-1\le \frac13(n-1)+\frac13\left(\frac12+\frac13+\cdots+\frac1n\right)\)，
所以
\(\frac1{a_{100}}\le 34+\frac13\left(\frac12+\frac13+\cdots+\frac1{100}\right) <34+\frac13\left(\frac12\times 6+\frac18\times 93\right)<40\)，
从而
\(100a_{100}>100\times \frac1{40}=\frac52\).
综上，
\(\frac52<100a_{100}<3\).
故选：B.
\end{solution}
\section{填空题}

\begin{problem}
我国南宋著名数学家秦九韶，发现了从三角形三边求面积的公式，他把这种方法称为“三斜求积”，它填补了我国传统数学的一个空白. 如果把这个方法写成公式，就是 \(S=\sqrt{\frac14\left[c^2a^2-\left(\frac{c^2+a^2-b^2}{2}\right)^2\right]}\)，其中 \(a,b,c\) 是三角形的三边，\(S\) 是三角形的面积. 设某三角形的三边 \(a=\sqrt2\)，\(b=\sqrt3\)，\(c=2\)，则该三角形的面积 \(S=\)\fillinblank{}.

\end{problem}

\begin{answer}
\(\dfrac{\sqrt{23}}{4}\)
\end{answer}

\begin{solution}
因为
\(S=\sqrt{\frac14\left[c^2a^2-\left(\frac{c^2+a^2-b^2}{2}\right)^2\right]}\)，
\(S=\sqrt{\frac14\left[4\times 2-\left(\frac{4+2-3}{2}\right)^2\right]}=\frac{\sqrt{23}}{4}\).
故答案为：\(\frac{\sqrt{23}}{4}\).
\end{solution}
\begin{problem}
已知多项式 \((x+2)(x-1)^4=a_0+a_1x+a_2x^2+a_3x^3+a_4x^4+a_5x^5\)，
则 \(a_2=\)\fillinblank{}，\(a_1+a_2+a_3+a_4+a_5=\)\fillinblank{}.

\end{problem}

\begin{answer}
\(8\)；\(-2\)
\end{answer}

\begin{solution}
含 \(x^2\) 的项为
\(x\cdot \binom43 x(-1)^3+2\cdot \binom42 x^2(-1)^2 =-4x^2+12x^2 =8x^2\)，
故 \(a_2=8\).

令 \(x=0\)，得 \(a_0=2\). 再令 \(x=1\)，得
\(0=a_0+a_1+a_2+a_3+a_4+a_5\)，
从而
\(a_1+a_2+a_3+a_4+a_5=-2\).
\end{solution}
\begin{problem}
若 \(3\sin\alpha-\sin\beta=\sqrt{10}\)，\(\alpha+\beta=\frac{\pi}{2}\)，
则 \(\sin\alpha=\)\fillinblank{}，\(\cos 2\beta=\)\fillinblank{}.

\end{problem}

\begin{answer}
\(\dfrac{3\sqrt{10}}{10}\)；\(\dfrac45\)
\end{answer}

\begin{solution}
方法一：利用辅助角公式处理.

因为 \(\alpha+\beta=\frac{\pi}{2}\)，所以 \(\sin\beta=\cos\alpha\)，即
\(3\sin\alpha-\cos\alpha=\sqrt{10}\).
即
\(\sqrt{10}\left(\frac{3\sqrt{10}}{10}\sin\alpha-\frac{\sqrt{10}}{10}\cos\alpha\right)=\sqrt{10}\).
令
\(\sin\theta=\frac{\sqrt{10}}{10}\)，\(\cos\theta=\frac{3\sqrt{10}}{10}\)，
则
\(\sqrt{10}\sin(\alpha-\theta)=\sqrt{10}\)，
所以
\(\alpha-\theta=\frac{\pi}{2}+2k\pi\)，\(k\in \Z\)，
即
\(\alpha=\theta+\frac{\pi}{2}+2k\pi\).
因此
\(\sin\alpha=\sin\left(\theta+\frac{\pi}{2}+2k\pi\right)=\cos\theta=\frac{3\sqrt{10}}{10}\).
则
\(\cos 2\beta=2\cos^2\beta-1=2\sin^2\alpha-1=\frac45\).

方法二：直接用同角三角函数关系式解方程.

因为 \(\alpha+\beta=\frac{\pi}{2}\)，所以 \(\sin\beta=\cos\alpha\)，即
\(3\sin\alpha-\cos\alpha=\sqrt{10}\).
又
\(\sin^2\alpha+\cos^2\alpha=1\)，
将 \(\cos\alpha=3\sin\alpha-\sqrt{10}\) 代入得
\(10\sin^2\alpha-6\sqrt{10}\sin\alpha+9=0\)，
解得
\(\sin\alpha=\frac{3\sqrt{10}}{10}\).
则
\(\cos 2\beta=2\cos^2\beta-1=2\sin^2\alpha-1=\frac45\).
\end{solution}
\begin{problem}
已知函数 \(f(x)=-x^2+2\)（\(x\le 1\)），\(f(x)=x+\dfrac{1}{x}-1\)（\(x>1\)）
则 \(f\!\left(f\!\left(\frac12\right)\right)=\)\fillinblank{}；若当 \(x\in[a,b]\) 时，\(1\le f(x)\le 3\)，则 \(b-a\) 的最大值是\fillinblank{}.

\end{problem}

\begin{answer}
\(\dfrac{37}{28}\)；\(3+\sqrt3\)
\end{answer}

\begin{solution}
先有
\(f\left(\frac12\right)=-\left(\frac12\right)^2+2=\frac74\)，
所以
\(f\!\left(f\!\left(\frac12\right)\right) =f\left(\frac74\right) =\frac74+\frac47-1 =\frac{37}{28}\).

当 \(x\le 1\) 时，由 \(1\le -x^2+2\le 3\) 得
\(-1\le x\le 1\).
当 \(x>1\) 时，由
\(1\le x+\frac1x-1\le 3\)
得
\(1<x\le 2+\sqrt3\).
因此
\([a,b]\subseteq [-1,\,2+\sqrt3]\)，
故
\(b-a\le (2+\sqrt3)-(-1)=3+\sqrt3\).
最大值为 \(3+\sqrt3\).
\end{solution}
\begin{problem}
现有 7 张卡片，分别写上数字 \(1,2,2,3,4,5,6\). 从这 7 张卡片中随机抽取 3 张，记所抽取卡片上数字的最小值为 \(\xi\)，则 \(P(\xi=2)=\fillinblank{}\)，\(E(\xi)=\fillinblank{}\).

\end{problem}

\begin{answer}
\(\dfrac{16}{35}\)；\(\dfrac{12}{7}\)
\end{answer}

\begin{solution}
从 7 张卡片中任取 3 张共有 \(\binom73\) 种方法.

最小值为 2 时，需从 \(\{2,2,3,4,5,6\}\) 中取 3 张且至少取到一个 2，其取法数为
\(\binom41+\binom12\binom42=16\)，
故
\(P(\xi=2)=\frac{16}{35}\).

\(\xi\) 的可能取值为 \(1,2,3,4\)，且
\(P(\xi=1)=\frac{15}{35}\)，\(P(\xi=2)=\frac{16}{35}\)，\(P(\xi=3)=\frac{3}{35}\)，\(P(\xi=4)=\frac{1}{35}\).
于是
\(E(\xi)=1\cdot \frac{15}{35}+2\cdot \frac{16}{35}+3\cdot \frac{3}{35}+4\cdot \frac{1}{35} =\frac{12}{7}\).
\end{solution}
\begin{problem}
已知双曲线 \(\frac{x^2}{a^2}-\frac{y^2}{b^2}=1\)（\(a>0\)，\(b>0\)） 的左焦点为 \(F\)，过 \(F\) 且斜率为 \(\frac{b}{4a}\) 的直线交双曲线于点 \(A(x_1,y_1)\)，交双曲线的渐近线于点 \(B(x_2,y_2)\)，且 \(x_1<0<x_2\). 若 \(|FB|=3|FA|\)，则双曲线的离心率是\fillinblank{}.

\end{problem}

\begin{answer}
\(\dfrac{3\sqrt6}{4}\)
\end{answer}

\begin{solution}
\begin{center}
\bitmapfigure[max width=0.26\textwidth,max totalheight=4.8cm]{img/2022/zhejiang/q16_solution_fig1.png}
\end{center}

设双曲线焦距参数为 \(c\)，则 \(F(-c,0)\).

过 \(F\) 的直线 \(AB\) 方程为
\(y=\frac{b}{4a}(x+c)\)，
右渐近线方程为
\(y=\frac{b}{a}x\).
联立得
\(B\left(\frac{c}{3},\frac{bc}{3a}\right)\).
又由 \(|FB|=3|FA|\) 且 \(A,B,F\) 共线，可得
\(A\left(-\frac{5c}{9},\frac{bc}{9a}\right)\).
将点 \(A\) 代入双曲线方程，
\(\frac{25c^2}{81a^2}-\frac{b^2c^2}{81a^2b^2}=1\)，
即
\(\frac{24c^2}{81a^2}=1\)，\(\frac{c^2}{a^2}=\frac{81}{24}\).
因此离心率
\(e=\frac{c}{a}=\sqrt{\frac{81}{24}}=\frac{3\sqrt6}{4}\).
\end{solution}
\begin{problem}
设点 \(P\) 在单位圆的内接正八边形 \(A_1A_2\cdots A_8\) 的边 \(A_1A_2\) 上，则 \(PA_1^2+PA_2^2+\cdots+PA_8^2\) 的取值范围是\fillinblank{}.

\end{problem}

\begin{answer}
\([12+2\sqrt2,\,16]\)
\end{answer}

\begin{solution}
\begin{center}
\bitmapfigure[max width=0.30\textwidth,max totalheight=4.8cm]{img/2022/zhejiang/q17_solution_fig1.png}
\end{center}

以圆心 \(O\) 为原点，建立平面直角坐标系，则
\(A_1(0,1),\ A_2\left(\frac{\sqrt2}{2},\frac{\sqrt2}{2}\right),\ \dots,\ A_8\left(-\frac{\sqrt2}{2},\frac{\sqrt2}{2}\right)\).
设 \(P(x,y)\)，则由对称性可得
\(PA_1^2+PA_2^2+\cdots+PA_8^2=8(x^2+y^2)+8\).
又因为 \(P\) 在线段 \(A_1A_2\) 上，所以
\(\cos 22.5^\circ\le |OP|\le 1\)，
即
\(\frac{1+\cos 45^\circ}{2}\le x^2+y^2\le 1\).
于是
\(8\cdot \frac{2+\sqrt2}{2}+8\le PA_1^2+\cdots+PA_8^2\le 16\)，
即
\([12+2\sqrt2,\,16]\).
\end{solution}
\section{解答题}

\begin{problem}
在 \(\triangle ABC\) 中，角 \(A,B,C\) 所对边分别为 \(a,b,c\). 已知 \(4a=\sqrt5\,c\)，\(\cos C=\frac35\).

\begin{enumerate}
\item 求 \(\sin A\) 的值；

\item 若 \(b=11\)，求 \(\triangle ABC\) 的面积.
\end{enumerate}
\end{problem}

\begin{solution}
（1）由 \(\cos C=\frac35\) 且 \(0<C<\pi\)，得
\(\sin C=\frac45\).
由正弦定理
\(\frac{a}{\sin A}=\frac{c}{\sin C}\)
及 \(4a=\sqrt5\,c\) 可得
\(4\sin A=\sqrt5\,\sin C\)，
故
\(\sin A=\frac{\sqrt5}{4}\cdot \frac45=\frac{\sqrt5}{5}\).

（2）由 \(4a=\sqrt5\,c\) 得 \(c=\frac{4a}{\sqrt5}\). 代入余弦定理
\(\cos C=\frac{a^2+b^2-c^2}{2ab}\)
并取 \(b=11\)，得
\(\frac{a^2+121-\frac{16}{5}a^2}{22a}=\frac35\).
化简得
\(a^2+6a-55=0\)，
解得 \(a=5\).

因此
\(S_{\triangle ABC}=\frac12ab\sin C =\frac12\cdot 5\cdot 11\cdot \frac45 =22\).
\end{solution}
\begin{problem}
如图，已知 \(ABCD\) 和 \(CDEF\) 都是直角梯形，\(AB\parallel DC\)，\(DC\parallel EF\)，\(AB=5\)，\(DC=3\)，\(EF=1\)，\(\angle BAD=\angle CDE=60^\circ\)，二面角 \(F-DC-B\) 的平面角为 \(60^\circ\). 设 \(M,N\) 分别为 \(AE,BC\) 的中点.



\begin{enumerate}
\item 证明：\(FN\perp AD\)；

\item 求直线 \(BM\) 与平面 \(ADE\) 所成角的正弦值.
\end{enumerate}

\bitmapfigure[max width=0.34\textwidth,max totalheight=4.8cm]{img/2022/zhejiang/q19_fig1.png}

\end{problem}

\begin{solution}
（1）过点 \(E,D\) 分别做直线 \(DC,AB\) 的垂线 \(EG,DH\)，并分别交于点 \(G,H\).

因为四边形 \(ABCD\) 和 \(EFCD\) 都是直角梯形，\(AB\parallel DC\)，\(CD\parallel EF\)，\(AB=5\)，\(DC=3\)，\(EF=1\)，\(\angle BAD=\angle CDE=60^\circ\)，由平面几何知识易知，
\(DG=AH=2\)，\(\angle EFC=\angle DCF=\angle DCB=\angle ABC=90^\circ\)，
则四边形 \(EFCG\) 和四边形 \(DCBH\) 是矩形，所以
\(EG=DH=2\sqrt3\)，\(FC=BC\).

又 \(DC\perp CF\)，\(DC\perp CB\)，且 \(CF\cap CB=C\)，所以 \(DC\perp\) 平面 \(BCF\). \(\angle BCF\) 是二面角 \(F-DC-B\) 的平面角，则
\(\angle BCF=60^\circ\).
所以 \(\triangle BCF\) 是正三角形. 由 \(DC\subset\) 平面 \(ABCD\)，得平面 \(ABCD\perp\) 平面 \(BCF\).

因为 \(N\) 是 \(BC\) 的中点，所以
\(FN\perp BC\).
又 \(DC\perp\) 平面 \(BCF\)，\(FN\subset\) 平面 \(BCF\)，可得
\(FN\perp CD\).
而 \(BC\cap CD=C\)，所以
\(FN\perp \text{平面 }ABCD\).
而 \(AD\subset\) 平面 \(ABCD\)，所以
\(FN\perp AD\).

（2）因为 \(FN\perp\) 平面 \(ABCD\)，过点 \(N\) 做 \(AB\) 平行线 \(NK\)，所以以点 \(N\) 为原点，\(NK,NB,NF\) 所在直线分别为 \(x\) 轴、\(y\) 轴、\(z\) 轴建立空间直角坐标系 \(N-xyz\).

\begin{center}
\bitmapfigure[max width=0.42\textwidth,max totalheight=5.2cm]{img/2022/zhejiang/q19_solution_fig1.png}
\end{center}

设
\(A(5,\sqrt3,0)\)，\(B(0,\sqrt3,0)\)，\(D(3,-\sqrt3,0)\)，\(E(1,0,3)\)，
则
\(M\left(3,\frac{\sqrt3}{2},\frac32\right)\)，
于是
\(\overrightarrow{BM}=\left(3,-\frac{\sqrt3}{2},\frac32\right)\)，\(\overrightarrow{AD}=(-2,-2\sqrt3,0)\)，\(\overrightarrow{DE}=(-2,\sqrt3,3)\).

设平面 \(ADE\) 的法向量为 \(\bs{n}=(x,y,z)\)，由
\(\bs{n}\cdot \overrightarrow{AD}=0\)，\(\bs{n}\cdot \overrightarrow{DE}=0\)
得
\[
\begin{cases}
-2x-2\sqrt3\,y=0,\\
-2x+\sqrt3\,y+3z=0,
\end{cases}
\]
取
\(\bs{n}=(\sqrt3,-1,\sqrt3)\).
设直线 \(BM\) 与平面 \(ADE\) 所成角为 \(\theta\)，则
\(\sin\theta=\frac{|\bs{n}\cdot \overrightarrow{BM}|}{|\bs{n}|\cdot |\overrightarrow{BM}|}=\frac{\left|3\sqrt3+\frac{\sqrt3}{2}+\frac{3\sqrt3}{2}\right|}{\sqrt{3+1+3}\cdot \sqrt{9+\frac34+\frac94}}=\frac{5\sqrt3}{\sqrt7\cdot 2\sqrt3}=\frac{5\sqrt7}{14}\).
\end{solution}
\begin{problem}
已知等差数列 \(\{a_n\}\) 的首项 \(a_1=-1\)，公差 \(d>1\). 记 \(\{a_n\}\) 的前 \(n\) 项和为 \(S_n\)（\(n\in\N^*\)）.

\begin{enumerate}
\item 若 \(S_4-2a_2a_3+6=0\)，求 \(S_n\)；

\item 若对于每个 \(n\in\N^*\)，存在实数 \(c_n\)，使 \(a_n+c_n,\ a_{n+1}+4c_n,\ a_{n+2}+15c_n\) 成等比数列，求 \(d\) 的取值范围.
\end{enumerate}
\end{problem}

\begin{solution}
（1）由
\(a_n=-1+(n-1)d\)
可得
\(S_4=4(-1)+6d=-4+6d\)，
且
\(a_2=-1+d\)，\(a_3=-1+2d\).
代入条件
\(S_4-2a_2a_3+6=0\)
得
\(-4+6d-2(-1+d)(-1+2d)+6=0\)，
化简得
\(d^2-3d=0\).
又 \(d>1\)，故 \(d=3\).

于是
\(a_n=3n-4\)，\(S_n=\frac{(a_1+a_n)n}{2}=\frac{3n^2-5n}{2}\).

（2）三项成等比数列，故
\((a_{n+1}+4c_n)^2=(a_n+c_n)(a_{n+2}+15c_n)\).
代入
\(a_n=nd-d-1\)
整理得关于 \(c_n\) 的二次方程
\(c_n^2+(14d-8nd+8)c_n+d^2=0\).
要使其有实数解，判别式需满足
\(\Delta=(14d-8nd+8)^2-4d^2\ge 0\).
化简为
\(\bigl[(n-2)d-1\bigr]\bigl[(2n-3)d-2\bigr]\ge 0\)（\(\forall n\in\N^*\)）.

取 \(n=2\)，得
\((-1)(d-2)\ge 0\)，
故 \(d\le 2\).

当 \(n\ge 3\) 时，由 \(d>1\) 可知上式恒成立. 因此
\(1<d\le 2\).
\end{solution}
\begin{problem}
如图，已知椭圆 \(\frac{x^2}{12}+y^2=1\). 设 \(A,B\) 是椭圆上异于 \(P(0,1)\) 的两点，且点 \(Q\left(0,\frac12\right)\) 在线段 \(AB\) 上，直线 \(PA,PB\) 分别交直线 \(y=-\frac12x+3\) 于 \(C,D\) 两点.



\begin{enumerate}
\item 求点 \(P\) 到椭圆上点的距离的最大值；

\item 求 \(|CD|\) 的最小值.
\end{enumerate}

\bitmapfigure[max width=0.40\textwidth,max totalheight=4.8cm]{img/2022/zhejiang/q21_fig1.png}

\end{problem}

\begin{solution}
（1）设椭圆上任意一点
\(H(2\sqrt3\cos\theta,\sin\theta)\).
则
\(|PH|^2 =12\cos^2\theta+(1-\sin\theta)^2 =13-11\sin^2\theta-2\sin\theta =-\!11\left(\sin\theta+\frac1{11}\right)^2+\frac{144}{11}\).
故
\(|PH|_{\max}=\sqrt{\frac{144}{11}}=\frac{12\sqrt{11}}{11}\).

（2）设直线 \(AB\) 的方程为
\(y=kx+\frac12\).
联立椭圆方程得
\(\left(k^2+\frac1{12}\right)x^2+kx-\frac34=0\).
设 \(A(x_1,y_1)\)，\(B(x_2,y_2)\)，则
\(x_1+x_2=-\frac{k}{k^2+\frac1{12}}\)，\(x_1x_2=-\frac{3}{4\left(k^2+\frac1{12}\right)}\).

再分别联立 \(PA,PB\) 与直线 \(y=-\frac12x+3\)，可化得
\(|CD|=\frac{3\sqrt5}{2}\cdot \frac{\sqrt{16k^2+1}}{|3k+1|}\).
由柯西不等式，
\(\frac{\sqrt{16k^2+1}}{|3k+1|} \ge \frac{4\cdot \frac34+1\cdot 1}{5} =\frac45\).
于是
\(|CD|\ge \frac{3\sqrt5}{2}\cdot \frac45=\frac{6\sqrt5}{5}\).
当 \(k=\frac{3}{16}\) 时取等号，故
\(|CD|_{\min}=\frac{6\sqrt5}{5}\).
\end{solution}
\begin{problem}
设函数 \(f(x)=\frac{e}{2x}+\ln x\)（\(x>0\)）.

\begin{enumerate}
\item 求 \(f(x)\) 的单调区间；

\item 已知 \(a,b\in\R\)，曲线 \(y=f(x)\) 上不同的三点 \((x_1,f(x_1))\)、\((x_2,f(x_2))\)、\((x_3,f(x_3))\) 处的切线都经过点 \((a,b)\). 证明：

\begin{enumerate}[label={（\arabic*）},itemsep=0.2em]
\item 若 \(a>e\)，则 \(0<b-f(a)<\frac12\left(\frac{a}{e}-1\right);\)
\item 若 \(0<a<e\)，\(x_1<x_2<x_3\)，则 \(\frac{2}{e}+\frac{e-a}{6e^2}<\frac1{x_1}+\frac1{x_3}<\frac{2}{a}-\frac{e-a}{6e^2}\).
\end{enumerate}
\end{enumerate}
\end{problem}

\begin{solution}
（1）函数导数为
\(f'(x)=-\frac{e}{2x^2}+\frac1x=\frac{2x-e}{2x^2}\).
故当 \(0<x<\frac{e}{2}\) 时，\(f'(x)<0\)；当 \(x>\frac{e}{2}\) 时，\(f'(x)>0\).

因此 \(f(x)\) 的减区间为
\(\left(0,\frac{e}{2}\right)\)，
增区间为
\(\left(\frac{e}{2},+\infty\right)\).

（2）设
\(g(x)=f(x)-b-f'(x)(x-a)\).
题设中三条不同切线都经过 \((a,b)\)，等价于方程 \(g(x)=0\) 有三个不同的正根 \(x_1,x_2,x_3\).

直接求导得
\(g'(x)=-\frac{1}{x^3}(x-e)(x-a)\).

（1）当 \(a>e\) 时，\(g\) 在 \((0,e)\)、\((a,+\infty)\) 上为减函数，在 \((e,a)\) 上为增函数. 因为 \(g\) 有 3 个不同的零点，故
\(g(e)<0<g(a)\).
由 \(g(e)<0\) 与 \(g(a)>0\) 分别化简得
\(b<\frac{a}{2e}+1\)，\(b>f(a)=\frac{e}{2a}+\ln a\).
于是
\(0<b-f(a)\)，
且
\(b-f(a)-\frac12\left(\frac{a}{e}-1\right) <\frac{a}{2e}+1-\left(\frac{e}{2a}+\ln a\right)-\frac12\left(\frac{a}{e}-1\right)\).
记右端为
\(u(a)=\frac32-\frac{e}{2a}-\ln a\).
则
\(u'(a)=\frac{e-2a}{2a^2}<0\)（\(a>e\)），
故 \(u(a)<u(e)=0\). 因此
\(0<b-f(a)<\frac12\left(\frac{a}{e}-1\right)\).

（2）当 \(0<a<e\) 时，同（1）中讨论可得：\(g\) 在 \((0,a)\)、\((e,+\infty)\) 上为减函数，在 \((a,e)\) 上为增函数.

不妨设 \(x_1<x_2<x_3\)，则
\(0<x_1<a<x_2<e<x_3\).
又因为 \(g(x)\) 有 3 个不同的零点，故
\(g(a)<0\)，\(g(e)>0\).
整理得到
\(\frac{a}{2e}+1<b<\frac{e}{2a}+\ln a\).

又
\(g(x)=1-\frac{a+e}{x}+\frac{ea}{2x^2}-\ln x+b\).
令
\(t=\frac{e}{x}\)，\(m=\frac{a}{e}\in (0,1)\)，
则方程
\(1-\frac{a+e}{x}+\frac{ea}{2x^2}-\ln x+b=0\)
即为
\(-(m+1)t+\frac{m}{2}t^2+\ln t+b=0\).

记
\(t_1=\frac{e}{x_1}\)，\(t_2=\frac{e}{x_2}\)，\(t_3=\frac{e}{x_3}\)，
则 \(t_1,t_2,t_3\) 为
\(-(m+1)t+\frac{m}{2}t^2+\ln t+b=0\)
的三个不同的根.

设
\(k=\frac{t_1}{t_3}=\frac{x_3}{x_1}>\frac{e}{a}>1\)，
则要证
\(\frac{2}{e}+\frac{e-a}{6e^2}<\frac1{x_1}+\frac1{x_3}<\frac{2}{a}-\frac{e-a}{6e^2}\)，
即证
\(2+\frac{e-a}{6e}<t_1+t_3<\frac{2e}{a}-\frac{e-a}{6e}\)，
即证
\(\frac{13-m}{6}<t_1+t_3<\frac{2}{m}-\frac{1-m}{6}\).
这又等价于
\(\left(t_1+t_3-\frac{13-m}{6}\right)\left(t_1+t_3-\frac{2}{m}+\frac{1-m}{6}\right)<0\)，
即证
\(t_1+t_3-2-\frac{2}{m}<\frac{(m-13)(m^2-m+12)}{36m(t_1+t_3)}\).

而
\(-(m+1)t_1+\frac{m}{2}t_1^2+\ln t_1+b=0\)，
且
\(-(m+1)t_3+\frac{m}{2}t_3^2+\ln t_3+b=0\)，
所以
\(\ln t_1-\ln t_3+\frac{m}{2}(t_1^2-t_3^2)-(m+1)(t_1-t_3)=0\)，
故
\(t_1+t_3-2-\frac{2}{m}=-\frac{2}{m}\times \frac{\ln t_1-\ln t_3}{t_1-t_3}\).
故即证
\(-\frac{2}{m}\times \frac{\ln t_1-\ln t_3}{t_1-t_3}<\frac{(m-13)(m^2-m+12)}{36m(t_1+t_3)}\)，
即证
\(\frac{(t_1+t_3)\ln \frac{t_1}{t_3}}{t_1-t_3}+\frac{(m-13)(m^2-m+12)}{72}>0\)，
即证
\(\frac{(k+1)\ln k}{k-1}+\frac{(m-13)(m^2-m+12)}{72}>0\).

记
\(\varphi(k)=\frac{(k+1)\ln k}{k-1}\)，\(k>1\)，
则
\(\varphi'(k)=\frac{1}{(k-1)^2}\left(k-\frac1k-2\ln k\right)\).
设
\(u(k)=k-\frac1k-2\ln k\)，
则
\(u'(k)=1+\frac1{k^2}-\frac2k>\frac2k-\frac2k=0\)，
所以
\(u(k)>u(1)=0\)，
从而
\(\varphi'(k)>0\).
故 \(\varphi(k)\) 在 \((1,+\infty)\) 上为增函数，所以
\(\varphi(k)>\varphi\left(\frac1m\right)\).

因此
\(\frac{(k+1)\ln k}{k-1}+\frac{(m-13)(m^2-m+12)}{72}> \frac{(m+1)\ln m}{m-1}+\frac{(m-13)(m^2-m+12)}{72}\).
记
\(\omega(m)=\ln m+\frac{(m-1)(m-13)(m^2-m+12)}{72(m+1)}\)，\(0<m<1\)，
则
\(\omega'(m)=\frac{(m-1)^2(3m^3-20m^2-49m+72)}{72m(m+1)^2} >\frac{(m-1)^2(3m^3+3)}{72m(m+1)^2}>0\).
所以 \(\omega(m)\) 在 \((0,1)\) 为增函数，故
\(\omega(m)<\omega(1)=0\)，
即
\(\ln m+\frac{(m-1)(m-13)(m^2-m+12)}{72(m+1)}<0\)，
从而
\(\frac{(m+1)\ln m}{m-1}+\frac{(m-13)(m^2-m+12)}{72}>0\).
故原不等式得证.
\end{solution}
